\section{Introduction}
\label{sec:intro}

Bound-To-Verdict (BTV)~\cite{btv2026} recently demonstrated that silent
decisions---AI-generated acts lacking a computationally bound evidence
chain---can be eradicated at the application layer. By encoding accountability
obligations as linear resource types in Safe Rust, BTV guarantees that no
\texttt{Verdict} object can materialise in memory without atomically consuming
its \texttt{EvidenceToken}. This shifts the enforcement of regulatory
obligations, such as LGPD Art.~18 and EU AI Act Art.~86, from runtime audit
checks to compile-time type constraints.

However, structural integrity in memory is not equivalent to verifiable
durability in the outside world. While BTV ensures that a decision \emph{was}
legitimate at the moment of construction, it does not mandate that the decision
\emph{remains} available for contestation. Existing solutions such as
Certificate Transparency (CT)~\cite{rfc6962} and Trillian~\cite{trillian}
provide robust append-only logs for certificates and arbitrary data, yet their
integration into AI pipelines remains voluntary---the decision to log is a
runtime policy choice, easily bypassed by an operator who controls the
infrastructure. An operator can satisfy the BTV invariant---generating a valid
\texttt{Verdict} in memory---while simply discarding the artefact before it
reaches persistent storage. We call this failure mode an \textbf{Ephemeral
Verdict} (Definition~\ref{def:ephemeral}): a decision that is legal in its
formation but inaccessible in its aftermath. The type \texttt{Verdict} proves
structural integrity; it does not prove persistence.

We introduce \textbf{BTV-Transparency}, a framework that extends the linear
type discipline to the persistence layer. We define a new type law for AI
accountability:
\begin{equation}
  \mathsf{DeliveryToken} \multimap (V \otimes R)
  \label{eq:delivery-law}
\end{equation}
where \texttt{DeliveryToken}---the sole type that authorises delivering a
decision to the end-user---can only be constructed by linearly consuming a
\texttt{Verdict}~($V$) and an \texttt{InclusionReceipt}~($R$). The
\texttt{InclusionReceipt} is a cryptographic proof, signed by a key the
operator does not control, asserting that the decision has been appended to a
public, verifiable transparency log.

This paper makes three contributions:
\begin{enumerate}
  \item \textbf{Linear Persistence Types (C1).} We define \texttt{InclusionReceipt}
  as a linear resource that encapsulates an Ed25519-signed Merkle-tree inclusion
  proof. We apply the same three protective barriers used in
  BTV---private fields, \texttt{pub(crate)} consumption, \texttt{\#[must\_use]}---to
  ensure that receipts cannot be forged or silently dropped.

  \item \textbf{The Persistence Enclosure Theorem (C2).} We prove that in any Safe
  Rust program integrating BTV-Transparency, a decision cannot be delivered to an
  end-user without a verified, non-repudiable proof of inclusion in an append-only
  log. The attempt to deliver an unpersisted decision is rejected by the compiler.

  \item \textbf{Empirical Validation (C3).} We provide a reference implementation in
  Safe Rust. Our evaluation demonstrates that enforcing verifiable persistence adds
  sub-millisecond overhead attributable exclusively to log-submission
  latency---negligible relative to the 50--200\,ms latency of typical LLM inference,
  proving that verifiable accountability is compatible with high-throughput systems.
\end{enumerate}

Together with BTV, this work closes the accountability loop: Paper~1 guarantees
that a decision without evidence cannot exist; Paper~2 guarantees that a decision
without verifiable persistence cannot be delivered.
